\numberlesssection{ВВЕДЕНИЕ}

Компьютерная графика играет ключевую роль в современном инженерном проектировании, образовании и моделировании физических процессов. Визуализация сложных механических систем, в том числе зубчатых передач, позволяет наглядно демонстрировать принципы их работы, выявлять ошибки проектирования на ранних этапах и облегчает понимание кинематических связей.

Статические чертежи и двухмерные схемы не всегда способны в полной мере передать динамику взаимодействия деталей, особенно когда речь идет о сложных цепочках зацепления. Построение интерактивных трехмерных моделей, позволяющих пользователю конструировать механизмы и наблюдать за их движением в реальном времени, является актуальной задачей.

\textbf{Цель курсовой работы} --- разработка программного обеспечения для интерактивного конструирования и визуализации кинематики шестереночных механизмов в трехмерном пространстве.

Для достижения поставленной цели необходимо решить следующие \textbf{задачи}:

\begin{itemize}
	\item выполнить анализ предметной области и существующих методов визуализации трехмерных сцен и моделирования кинематических связей;
	\item выбрать алгоритмы компьютерной графики для визуализации сцены: алгоритмы удаления невидимых поверхностей, методы закраски и построения теней;
	\item разработать выбранные алгоритмы визуализации;
	\item разработать алгоритм расчета кинематики зубчатых передач для корректного вычисления скоростей и направлений вращения всех элементов цепи;
	\item спроектировать структуру программного обеспечения и реализовать его;
	\item провести тестирование программы и оценить эффективность реализованных алгоритмов.
\end{itemize}
